% !TeX root = ../papers/nc-p3.tex
\section{二重次数の全境界と第二解の閉包}
\label{sec:klein-doublet-boundary}
\begin{theorem}\label{thm:klein-doublet-complete-cover}
次数次元$(2,2,1,1)$の片側対称な標準AS窓は、両対称チャートまたは
定理\ref{thm:klein-two-doublets-generic}の第二解の閉包に属する。
式\eqref{eq:doublet-completions}の$\det X=0$の全入力では両対称チャートの
判定で正則性が決まる。$\det X\ne0$では薄い次数の成分が零の場合を含め、
相異なる全ての行列補完が同じ既約族の閉包に入る。
後者の個々の特殊化での正則性は別に判定する必要がある。
\end{theorem}

\paragraph{交差ブロックの階数が落ちる場合。}
まず薄い次数の全成分が非零なら、前節の倍率変更で$U,V$の全成分を1にできる。
実際、$\det X=0$は二つの積比が等しい条件である。
$X,Y$は全成分1の行列$J$になる。$h=0,1$に対し、$M_h$の列空間の
$A,B$への射影階数を$u_{Ah},u_{Bh}$、行空間について$v_{Ah},v_{Bh}$とする。
各階数は1または2で、階数1なら像は$(1,1)$の直線である。
$u_{Ah}=v_{Bh}=2$は同時に起きない。階数2の因子分解を$A$行、$B$列に
制限すると可逆な二行列の積になり、交差ブロックの階数1に反する。
$u_{Bh},v_{Ah}$についても同様である。

$u_{A0}=u_{A1}=1$なら、$x=e_0+e_1$で偶評価すると、厚い二次数からの
奇$A$上の二式は共に$(1,1)$の倍数であり、薄い次数からの奇$B$上の
二式も比例する。評価の総階数は2以下である。他の組と相手窓についても同様である。
従って$u_{A0},u_{A1}$のうち一方だけが2であり、$v_{B0},v_{B1}$では
他方だけが2である。もう一つの組合せ$u_B,v_A$もそうなる。
もし$u_A,u_B$の階数2の色が同じなら、色1で両組の列射影がともに1か、
両組の行射影がともに1である。該当する窓の$01,10$の係数は同じで、
$23,32$も同じになる。色0はもとから対角、薄い次数の係数は全て1なので、
この窓は隣接対称である。二つの正交換から標準ASでは完全対称となる。

色が反対なら、たとえば色0で$A$の行・列射影がともに2、$B$がともに1である。
色1では逆になる。従って$T,S'$が$J$の定数倍であり、$S,S'$および
$T,T'$の非対角成分の一致から$S,T'$も対称行列である。組や色を逆にしても同じである。
$M_0,M_1$と薄い二行列は対称で、テンソルには二文字ずつの巡回交換がある。
これと偶交換を合成すると奇交換が得られる。

次に薄い成分に零を持つ場合、各行・列で零は高々一つである。$U_{00}=0$としてよい
（両組や組内のラベルを変え、必要なら薄い二色を交換する）。
$\det X=0$から$V_{00}V_{11}=0$である。
同じ薄い色の四成分中三つが零なら、残る一成分が純関係を与えるか隣接階数が落ちる。
従って両行列に丁度一つずつ、同じ色に零がある。$V_{11}=0$なら
$\gamma$の二行列で列を反転した零点同士が転置位置にあるため、
窓を取り替え、組内を表示し直せばよい。従って
\[
 U=\begin{pmatrix}0&u\\v&w\end{pmatrix},\quad
 V=\begin{pmatrix}0&p\\q&s\end{pmatrix},\quad
 X=\begin{pmatrix}0&0\\pw&sv\end{pmatrix},\quad
 Y=\begin{pmatrix}0&0\\us&wq\end{pmatrix}
\]
として六係数は非零である。$M_1^{\mathsf t}$の両組への列射影はともに2である。
もし1なら、交差ブロックの列方向が$(0,1)$なので一行全部が零になり、
厚い次数の座標評価成分が零になる。ゆえに$M_1$の両組の列射影はともに1で、
その二方向は$(us,wq),(pw,sv)$である。
奇座標を$f_{B1}/f_{A0}=q/u$, $f_{B0}/f_{A1}=p/v$,
$f_{B1}/f_{A1}=s/w$の倍率比で変える（ここだけ$f$は倍率）。
これは薄い二行列を転置位置で対称にし、同時に
$f_{A1}/f_{A0}=wq/(us)$, $f_{B1}/f_{B0}=sv/(pw)$により厚い色1でも
隣接対称にする。従ってこの場合の正則入力も完全対称である。

\paragraph{可逆な交差ブロックと全ての補完。}
$h=\det S'/\det S\ne0$とする。$h\ne1$では
$Z=(S'-hS)/(1-h)$と置くと、$S=Z+\operatorname{diag}(v,w)$,
$S'=Z+h\operatorname{diag}(v,w)$である。二行二列の余因子行列は線型なので、
下側の非対角成分の一致は
$Y\operatorname{adj}(Z)X=\operatorname{diag}(s,t)$と同値となる。
従って$Z$は$s,t$から一意に決まり、$\det Z=st/(\det X\det Y)$である。
二つの行列式の比を取る残りの条件はただ一つ
\[
             h v w=\det Z
\]
となる。この表示は可逆性を課す開部分で逆写像も持つ。基底$U,V$の八係数と
$h$について$h\det X\det Y$は単元であり、この方程式は$vw$と$st$を
結ぶ既約な二項式である（$v$について原始な一次多項式）。
空でない開条件を取り除いても既約であり、一般の薄い全成分非零の第二解が稠密である。

$h=1,S'\ne S$では$S'-S=\operatorname{diag}(m,n)$で、
$Y\operatorname{adj}(S'-S)X$の非対角が零になる。
$U=\left(\begin{smallmatrix}u&v\\w&d\end{smallmatrix}\right)$,
$V=\left(\begin{smallmatrix}p&q\\r&s\end{smallmatrix}\right)$と書くと、
$(m,n)$にかかる二行列の行列式は
$(psvw-rqud)(psvw+rqud)$となる。第一因子が$\det X\ne0$だから
第二が零で、八係数は全て非零である。前節の規格化で$U=J$,
$V=\left(\begin{smallmatrix}1&1\\1&z\end{smallmatrix}\right)$とすれば$z=-1$である。
$S=\left(\begin{smallmatrix}a&b\\c&d\end{smallmatrix}\right)$と書き直す。
このとき非対角条件は$m=n$で、行列式を保つ非自明解は
$(a',d')=(-d,-a)$で尽きる。

これも同じ閉包に入ることを直接示せる。$z\ne\pm1$で$k=(zb+c)/(z+1)$とすれば
$Y\operatorname{adj}(Z)X$を対角にする中心行列$Z$は
$Z_{00}=Z_{11}=k$, $Z_{01}=b,Z_{10}=c$である。従って第二解は
\[
 a'=k+\frac{k^2-bc}{d-k},\qquad d'=k+\frac{k^2-bc}{a-k}.
\]
$b\ne c$なら$z\to-1$で$k$は極を持ち、二値は$-d,-a$に延びる。
$b=c$の同じ解は$z=-1$で$c$を$b$に近づけて得る。
これで定理\ref{thm:klein-doublet-complete-cover}が従う。

\begin{remark}
最小周期4の主枝で正規パラメータを探すなら、次数にも制約がある。
各位置にある正則な正規元の族$q_i\in A_{i,i+n}$が
$A_{ij}q_j=q_i A_{i+n,j+n}$を全ての$i,j$で満たせば、
$a q_j=q_i\sigma(a)$で一意に定まる$\sigma$は乗法的な線型全単射である。
従って$n$は周期であり4の倍数になる。特にこの枝で先の二周期窓の
正則二次パラメータによる証明をそのまま適用することはできない。
\end{remark}

\begin{theorem}[有理閉点での最小周期4の正則性]\label{thm:klein-rational-period-four}
節\ref{sec:klein-two-doublets}で二十点の消去に使った有理入力自身が、
最小周期4、反対階数10の標準Koszul AS正則窓である。
四頂点で折り畳んだ環は四つの中心四次式の多項式環上で有限自由であり、
各始点での自由階数は256である。
\end{theorem}
\begin{proof}
\path{computations/verify_klein_period_four_center.py}で次の有限証明書を再生成する。
各内部次数・文字次数で左端に一文字を付け、先頭二文字の関係に既知の商基底を
右から掛けて掃き出す。次数5までの元の次元は各始点で$1,4,10,20,35,56$である。
各位置の四次式$q_p$の係数に$ q_p v_i=v_i q_{p+1}$を課す。添字$p$は4を法とする。
有理数体上の同次解空間の次元は、四つの文字次数順に$(3,2,0,0)$である。
計算コードに指定した有理核基底について、次数零の側から係数$(1,3,5),(2,7,11)$の二式、
次数1の側から基底二式を取る。各四式の係数の分母を払い原始整数化する。
これらは折り畳み環の中心に入る。

四式を関係に追加し、各始点で二次関係と四次関係の両方に右から商の基底語を掛けて
同じ次数別掃き出しを行う。素数1009で得る商の次元は次数0から13まで
\[
 1,4,10,20,31,40,44,40,31,20,10,4,1,0.
\]
従って標数零での各次元にはこれが上界となり、一次生成性で以後は零である。
この上界の母関数は$(1+t+t^2+t^3)^4$である。中心性により商の基底を持ち上げると、
重み4の四変数多項式環上の生成系を得る。その自由加群による上界は$(1-t)^{-4}$となる。
一方、当該有理点は正則な幾何学的一般点を持つ第二解の族に属し、各次数で関係を
挿入する行列の余核次元は特殊化で減少しない。従って全次数で上下界が等しく、
四式は代数的に独立で、この生成系は自由基底である。

四階テンソルの一文字と隣接の切断はそれぞれ全階数4と6である。
四頂点で折り畳んだ右Koszul列の最初の二微分は表示から完全で、最後の二微分は
テンソルの完全対合により左の表示の双対である。
多項式環上で自由なので、始点単純から環への$\operatorname{Hom}$と
$\operatorname{Ext}^1$は零となる。実際、左増大イデアル上の写像$f$に対し、
中心変数のうち二つを$u,v$として$u f(v)=v f(u)$を使えば、
互いに素な変数であることから$f(u)=u b$となり、$u f(z)=z f(u)$で
$f(z)=zb$が従う。よって双対側も上の二項で完全で、中央はHilbert級数の
Euler等式から完全になる。左右を逆にしても同様である。
これで全次数のKoszul性と左右AS性を得る。構成した有理点で三配置の非同値性は
節\ref{sec:klein-generic-isomorphisms}の同じ素数12619での比較から直接従い、
反対階数と最小周期も主張の値になる。
\end{proof}
